\chapter{Introduction}\label{ch:introduction}

\section{The problem}\label{sec:problem}

A multicopter is unstable without feedback control.
The flight controller must hold the attitude, follow position commands and reject wind, all at once and on different airframes.
PX4 does this with a cascade of proportional and PID controllers.
Their gains must be tuned for each airframe.
The tuning is valid for one mass, one motor and propeller set and one battery state.
When the vehicle changes, for example a heavier payload or a weaker battery, the tuned gains no longer match.

Incremental nonlinear dynamic inversion (INDI) offers another route.
It does not model the whole vehicle.
It measures how the vehicle actually accelerates and asks the actuators for the change that closes the remaining gap.
The method has been flown on aircraft, micro air vehicles and quadrotors in aggressive flight~\cite{sieberling2010,smeur2016,smeur2018,tal2021}.
This report describes an implementation inside PX4 that works in every multicopter flight mode and tests how far it can be trusted.

\section{Why INDI}\label{sec:why}

The idea is explained in \cref{ch:background}.
Here we list what the designer gains.
\begin{itemize}
  \item \textbf{Few model parameters.}
    The inner loop needs the control effectiveness of each axis and the time constant of the motors.
    It needs no drag model and no model of the other forces and torques on the vehicle.
  \item \textbf{Disturbance rejection without an observer.}
    A disturbance shows up in the measured acceleration.
    The next command corrects it, so no integrator and no disturbance estimator are required.
  \item \textbf{One gain per axis.}
    The loop bandwidth of each rate axis is set by one gain.
  \item \textbf{Use of the rotor speed.}
    Bidirectional DShot telemetry reports the speed of each rotor.
    The controller knows what the motors really do, not only what was commanded.
\end{itemize}
There is a price.
The law needs a clean, fast measurement of angular and linear acceleration, and it relies on the actuator model being close.
Later chapters show where this price appears.

\section{Goals}\label{sec:goals}

The project has four goals.
\begin{enumerate}
  \item Build an INDI angular-rate controller and an INDI linear-acceleration controller for PX4.
  \item Derive them from first principles and analyse their stability margins.
  \item Test them against acceptance criteria that were written before the tests.
  \item State clearly what does not work yet and what must be checked on real hardware.
\end{enumerate}

\section{Scope and constraints}\label{sec:scope}

The following rules shaped the design.
\begin{description}
  \item[PID stays untouched and default.] The stock controllers are not modified.
    The INDI modules are selected by two parameters, \param{MC_INDI_EN} for the rate loop and \param{MC_INDI_POS_EN} for the acceleration loop.
    With both at zero the vehicle flies exactly as before.
  \item[Every flight mode.] Manual, assisted, automatic and offboard modes all go through the same two modules.
    There is no fall-back to the PID controller inside them.
  \item[Generic design, one reference airframe.] The design makes no assumption specific to one vehicle.
    The tests use one simulated 5-inch quadcopter of about 600~g (\cref{sec:airframe}).
    Defaults for other vehicles are a stated limit.
  \item[Simulation only.] All evidence comes from Gazebo software-in-the-loop (SITL) runs.
    No real flight was made.
    \Cref{ch:limits} and \cref{ch:bringup} say what this means.
\end{description}

\section{Acceptance criteria and results}\label{sec:criteria}

Eight acceptance criteria were fixed before the tests.
\Cref{tab:criteria} lists them with their verdicts.
\Cref{ch:validation} gives the data, and \cref{ch:traceability} maps each criterion to its evidence.

\begin{table}[tbp]
\caption{The eight acceptance criteria and their verdicts.}\label{tab:criteria}
\centering
\small
\begin{tabularx}{\textwidth}{@{}l l Y l@{}}
\toprule
 & Criterion & Meaning & Verdict \\
\midrule
C1 & Evidence integrity & The logs and scoring scripts can be trusted: complete logs, two independent scorers, recorded provenance of every run. & \passed \\
C2 & Stability margins & Phase margin of at least 45 deg and gain margin of at least 6 dB, measured in flight by multisine injection, and clean attitude steps. & \passed \\
C3 & Robustness & A campaign of 300 randomised simulated flights: no loss of control and no sustained oscillation. & \notpassed \\
C4 & Noise and effort & Under three noise profiles, motor effort and command chatter stay within limits. & \passed \\
C5 & Modes and transitions & All flight modes, mode switches, failsafes, kill and re-arm behave as specified. & \passed \\
C6 & Defaults and tuning method & Starting from wrong physical parameters, the identification chain must lead to a controller that passes C2, C3 and C5. It requires C3. & \notpassed \\
C7 & Code and upstream readiness & Stock modules untouched, unit tests, code format, flash size and review. & \passed \\
C8 & Report & Design, analysis, evidence, limits and bring-up are written down. & \passed \\
\bottomrule
\end{tabularx}
\end{table}

\begin{keyresult}[Result in one paragraph]
Six of eight criteria are met.
C3 and C6 are not.
The controllers are stable with good margins and handle noise, modes and transitions in simulation.
They are not robust enough over the full randomised campaign.
Two failures belong to INDI alone: PID flown on the same draws does not fail them.
\end{keyresult}

\section{Reading guide}\label{sec:guide}

The chapters follow one line: problem, idea, design, analysis, validation, limits, bring-up.
\begin{itemize}
  \item \Cref{ch:background} explains the INDI principle, the PX4 control chain, the reference airframe and the notation.
  \item \Cref{ch:inner_loop} and \cref{ch:outer_loop} give the design of the two controllers.
  \item \Cref{ch:analysis} analyses stability and robustness.
  \item \Cref{ch:validation} reports the simulation evidence against the criteria.
  \item \Cref{ch:limits} states what does not work yet.
  \item \Cref{ch:bringup} lists the checks needed before a first real flight.
  \item \Cref{ch:conclusion} concludes.
  \item The appendices list the parameters and map each criterion to its evidence.
\end{itemize}
A reader who knows control theory but not this project should read \cref{ch:background} first, then the chapters in order.
A reader who only wants the verdicts can read \cref{sec:criteria} and \cref{ch:limits}.
